\section{Theoretical Background}
\label{sec:theory}

\subsection{Signal Chain}
\label{sec:signal-chain}

The complete signal chain from the digital-to-analog converter (DAC) output
voltage $\Vda$ to the Hall sensor output voltage $\Vm$ is shown in
\cref{fig:signal-chain}. The chain consists of five stages:

\begin{enumerate}
\item \textbf{Amplifier}: The DAC voltage $\Vda$ is converted to a coil current
$I = \kA \cdot \Vda$, where $\kA = \SI{0.3614}{\ampere/\volt}$ is the
transconductance amplifier gain.

\item \textbf{Magnetic circuit}: The current $I$ flowing through a coil of
$\Nc = 50$ turns produces a magnetic flux $\Phi = \kpole \cdot I$, where
$\kpole = \Nc / \Ra$ is the flux-per-ampere constant and $\Ra$ is the total
reluctance of the magnetic circuit.

\item \textbf{Monopole model}: The flux is modeled as originating from an
equivalent magnetic charge $q = \Phi / \mu_0$ located at a distance $b$
behind the pole tip surface.

\item \textbf{Field decay}: The magnetic field at a distance $x$ from the
pole tip is $B(x) = \kpole \cdot I / [4\pi (x+b)^2]$, where $\mu_0$
cancels between the charge definition and Coulomb's law.

\item \textbf{Hall sensor}: The EQ-730L sensor~\cite{EQ730L} converts the
field to a voltage $\Vm = \SH \cdot B$, with sensitivity
$\SH = \SI{130}{\volt/\tesla}$.
\end{enumerate}

% ---- TikZ: Signal chain block diagram (NEW-1) ----
\begin{figure}[htbp]
\centering
\begin{tikzpicture}[
    block/.style={draw, thick, minimum width=2.2cm, minimum height=1.0cm,
                  font=\small},
    arr/.style={-{Stealth[length=3mm]}, thick},
    lbl/.style={font=\footnotesize, midway, above}
]
% Blocks
\node[block] (dac) {DAC};
\node[block, right=1.5cm of dac] (amp) {Amplifier};
\node[block, right=1.5cm of amp] (coil) {Coil};
\node[block, right=1.5cm of coil] (mono) {Monopole};
\node[block, right=1.5cm of mono] (hall) {Hall Sensor};

% Arrows with signal labels
\draw[arr] (dac) -- node[lbl] {$\Vda$} (amp);
\draw[arr] (amp) -- node[lbl] {$I = \kA \Vda$} (coil);
\draw[arr] (coil) -- node[lbl] {$\Phi = \kpole I$} (mono);
\draw[arr] (mono) -- node[lbl] {$B(x)$} (hall);

% Output
\draw[arr] (hall.east) -- ++(1.2,0) node[right, font=\small] {$\Vm = \SH \, B$};

% Gain labels below blocks
\node[font=\scriptsize, below=0.15cm of amp] {$\kA$};
\node[font=\scriptsize, below=0.15cm of coil] {$\Nc,\;\Ra$};
\node[font=\scriptsize, below=0.15cm of mono] {$q = \Phi/\mu_0$};
\node[font=\scriptsize, below=0.15cm of hall] {$\SH$};
\end{tikzpicture}
\caption{Signal chain block diagram. The DAC voltage $\Vda$ drives a
transconductance amplifier ($\kA$), which produces a coil current $I$.
The magnetic circuit converts $I$ to flux $\Phi$ and then to field $B(x)$
via the monopole model. The Hall sensor outputs $\Vm \propto B$.}
\label{fig:signal-chain}
\end{figure}

\subsection{Monopole Model}
\label{sec:monopole-model}

We model each pole tip as a point magnetic charge $q$ located at a distance
$b$ behind the physical tip surface (\cref{fig:monopole-geometry}). This is
the standard approach used by Zhang and Menq~\cite{Zhang2010,Menq2011}.

The magnetic flux produced by a single pole carrying current $I$ is
\begin{equation}
\Phi = \kpole \cdot I = \frac{\Nc}{\Ra} \cdot I,
\label{eq:flux}
\end{equation}
where $\kpole \equiv \Nc / \Ra$ has units of Wb/A. The equivalent magnetic
charge is
\begin{equation}
q = \frac{\Phi}{\mu_0} = \frac{\kpole \cdot I}{\mu_0}.
\label{eq:charge}
\end{equation}

Applying Coulomb's law for magnetic charges, the field at distance $r = x + b$
from the charge is
\begin{equation}
B(x) = \frac{\mu_0}{4\pi} \cdot \frac{q}{(x+b)^2}
     = \frac{\mu_0}{4\pi} \cdot \frac{\kpole \cdot I}{\mu_0 (x+b)^2}
     = \frac{\kpole \cdot I}{4\pi (x+b)^2}.
\label{eq:B-field}
\end{equation}
Crucially, $\mu_0$ cancels between the charge definition~\eqref{eq:charge}
and Coulomb's law, so $B$ depends only on $\kpole$, $I$, and the geometry.

% ---- TikZ: Monopole geometry (NEW-2) ----
\begin{figure}[htbp]
\centering
\begin{tikzpicture}[>=Stealth, thick]
% Pole body (shaded rectangle)
\fill[gray!20] (-3.5,-.8) rectangle (-0.3,.8);
\draw[thick] (-3.5,-.8) rectangle (-0.3,.8);
\node at (-1.9, 0) {Pole};

% Pole tip (rounded)
\draw[thick] (-0.3, 0.5) .. controls (0.1, 0.5) and (0.2, 0.3) .. (0.2, 0)
             .. controls (0.2, -0.3) and (0.1, -0.5) .. (-0.3, -0.5);

% Tip surface line
\draw[dashed, gray] (0.2, -1.2) -- (0.2, 1.2) node[above, font=\small, black] {tip surface};

% Charge location (behind tip)
\node[circle, fill=black, inner sep=2pt, label={below:$q$}] (charge) at (-1.0, 0) {};

% Distance b (charge to tip)
\draw[{Stealth}-{Stealth}, thick] (-1.0, -1.0) -- (0.2, -1.0);
\node[below, font=\small] at (-0.4, -1.0) {$b$};

% Hall sensor
\node[draw, thick, minimum width=0.6cm, minimum height=1.0cm,
      fill=blue!10] (sensor) at (3.5, 0) {};
\node[font=\footnotesize, below=0.1cm of sensor] {Hall};

% Distance x (tip to sensor)
\draw[{Stealth}-{Stealth}, thick] (0.2, 1.0) -- (3.2, 1.0);
\node[above, font=\small] at (1.7, 1.0) {$x$};

% Total distance r = x + b
\draw[{Stealth}-{Stealth}, thick] (-1.0, -1.6) -- (3.2, -1.6);
\node[below, font=\small] at (1.1, -1.6) {$r = x + b$};

% Field arrow
\draw[-{Stealth}, thick, blue] (0.5, 0) -- (3.0, 0)
    node[midway, above, font=\small, blue] {$B(x)$};
\end{tikzpicture}
\caption{Monopole geometry. The equivalent charge $q$ resides at distance
$b$ behind the tip surface. The Hall sensor at distance $x$ from the tip
measures $B(x) = \kpole I / [4\pi(x+b)^2]$.}
\label{fig:monopole-geometry}
\end{figure}

\subsection{Transfer Function and Linearization}
\label{sec:linearization}

Combining the signal chain, the measured transfer function (ratio of Hall
sensor output to DAC input) at distance $x$ is
\begin{equation}
H(x) = \frac{\Vm}{\Vda}
     = \frac{\SH \cdot \kpole \cdot \kA}{4\pi} \cdot \frac{1}{(x+b)^2}
     \equiv \frac{a^2}{(x+b)^2},
\label{eq:H-model}
\end{equation}
where
\begin{equation}
a^2 = \frac{\SH \cdot \kpole \cdot \kA}{4\pi} \times 10^{12}
\label{eq:a-squared}
\end{equation}
when $x$ and $b$ are expressed in micrometers ($10^{12}$ converts
$\text{m}^{-2}$ to $\mu\text{m}^{-2}$).

For parameter estimation, we linearize~\eqref{eq:H-model} by taking the
reciprocal square root:
\begin{equation}
\frac{1}{\sqrt{H(x)}} = \frac{x + b}{a} = \frac{1}{a} \cdot x + \frac{b}{a}.
\label{eq:linearized}
\end{equation}
This is a linear function of $x$ with slope $1/a$ and intercept $b/a$.
Ordinary least-squares regression yields both $a$ and $b$ directly, with
$R^2$ computed in the linearized space.

\subsection{Parameter Recovery}
\label{sec:parameter-recovery}

From the fitted value of $a$, the physical parameters are recovered as:
\begin{align}
\kpole &= \frac{4\pi \, a^2}{\SH \cdot \kA \times 10^{12}},
\label{eq:kpole} \\
\Ra &= \frac{\Nc}{\kpole}.
\label{eq:Ra}
\end{align}

The magnetic force on a paramagnetic bead of volume $\Vbead$ and effective
susceptibility $\chitot$ is
\begin{equation}
\Fmag(x) = \frac{\chitot \, \Vbead}{\mu_0} \, B(x) \left|\frac{\mathrm{d}B}{\mathrm{d}x}\right|,
\label{eq:force}
\end{equation}
where the field gradient is
\begin{equation}
\frac{\mathrm{d}B}{\mathrm{d}x} = -\frac{2 \, \kpole \cdot I}{4\pi (x+b)^3}.
\label{eq:gradient}
\end{equation}
